% Lean module: MonadFormal.Information.Testing.MovingSupport
\subsubsection{A moving-support counterexample}
\label{sec:moving-support}

Let $U_\theta$ be the Lebesgue-uniform probability law on $(0,\theta]$,
where $\theta>0$. Thus
$U_\theta=\theta^{-1}\lambda|_{(0,\theta]}$, with $\lambda$
Lebesgue measure. The endpoint is an unknown parameter, so changing it
changes the support, not just the density on a fixed set.
\leandef{MonadFormal.Information.intervalLaw}

\begin{proposition}[Exact endpoint-testing error]
\label{prop:uniform-testing-error}
For $0<\theta\le u$, arbitrary binary priors $a,b$, and $n\ge0$,
\[
 e_\pi(U_\theta^{\otimes n},U_u^{\otimes n})
 =\min\left\{a,b\left(\frac\theta u\right)^n\right\}.
\]
For balanced priors this is
$\tfrac12(\theta/u)^n$.
\end{proposition}
\begin{proof}
Take $A=(0,\theta]$. Direct restriction of Lebesgue measure gives
$U_\theta(A)=1$, $U_u(A)=\theta/u$ and
$U_u|_A=(\theta/u)U_\theta$.
Apply \Cref{cor:conditioning-iid}. Since $0<\theta/u\le1$,
the second term in the minimum is no larger than the first when
$a=b=1/2$.
\leanproof{MonadFormal.Information.intervalLaw_mass,MonadFormal.Information.intervalLaw_restrict,MonadFormal.Information.intervalLaw_binaryError,MonadFormal.Information.intervalLaw_binaryError_toReal,MonadFormal.Information.intervalLaw_binaryError_equalPrior}
\end{proof}

\begin{theorem}[Inverse-sample local scale]
\label{thm:uniform-local-scale}
Fix $\theta>0$ and $h\ge0$, and for $n\ge1$ put
$u_n=\theta+h/n$. Then
\[
 \lim_{n\to\infty}e_\pi(U_\theta^{\otimes n},U_{u_n}^{\otimes n})
 =\min\{a,b e^{-h/\theta}\}.
\]
With balanced priors and $h>0$, this limit is
$\tfrac12 e^{-h/\theta}\in(0,1/2)$.
\end{theorem}
\begin{proof}
The preceding proposition reduces the limit to
\[
 \left(\frac{\theta}{\theta+h/n}\right)^n
 =\left(1+\frac{h/\theta}{n}\right)^{-n}
 \longrightarrow e^{-h/\theta}.
\]
The usual exponential limit is positive, so taking reciprocals
preserves convergence. Multiplication by $b$ and taking the minimum
with $a$ are continuous. For $h>0$, $0<e^{-h/\theta}<1$,
which gives the balanced-prior assertion.
\leanproof{MonadFormal.Information.tendsto_interval_support_power,MonadFormal.Information.tendsto_interval_binaryError,MonadFormal.Information.tendsto_interval_binaryError_equalPrior,MonadFormal.Information.interval_local_error_nontrivial}
\end{proof}

The alternative changes with $n$: this is a local testing limit,
not the fixed-alternative Chernoff exponent of
\Cref{thm:chernoff-exponent}. A displacement of order $n^{-1}$
already gives a nontrivial testing problem here. A universal
$n^{-1/2}$ parameter scale would therefore erase the dependence on
the family's regularity and support behavior.
